\begingroup
\renewcommand{\arraystretch}{1.45}
\fancyhead[L]{\small ARITHMETIC–GEOMETRIC MEAN.}
\Needspace{34\baselineskip}
\AGMArticle{13}
[The first-order differentials of the individual terms of the complete algorithm for an arithmetic–geometric mean are related to one another in two ways in Article 9. One relation involves only the differentials of the logarithms of the quotients of those quantities. It follows from repeated application of the two relations $aa=bb+cc$, $\frac{a'}{c'}=\frac{a+b}{a-b}$, and, by using the results obtained in the preceding article, can be extended to determine the differential of the quotient of two corresponding arithmetic–geometric means and expressed so that $\frac{1}{cc}\,\mathrm{d}\log\frac{a}{b}$, which may be denoted by $\Delta$, has the following expressions for every positive and negative $n$, where a negative following index $n$ is interpreted as equivalent to a positive preceding index of the same absolute value:
\SourceLeaf{389}
\Needspace{10\baselineskip}
\[
\begin{aligned}&\\&=\frac{1}{2^{n}}\mathbin{.}\frac{1}{a^{\mathrm{n}}a^{\mathrm{n}}}\,\mathrm{d}\log\frac{c^{\mathrm{n}}}{b^{\mathrm{n}}}\\&=\frac{1}{2^{n}}\mathbin{.}\frac{1}{b^{\mathrm{n}}b^{\mathrm{n}}}\,\mathrm{d}\log\frac{c^{\mathrm{n}}}{a^{\mathrm{n}}}\\&=\frac{1}{2^{n}}\mathbin{.}\frac{1}{c^{\mathrm{n}}c^{\mathrm{n}}}\,\mathrm{d}\log\frac{a^{\mathrm{n}}}{b^{\mathrm{n}}}\\&=\frac{\pi}{2}\mathbin{.}\frac{1}{\mathrm{M}(a,b)\mathbin{.}\mathrm{M}(a,c)}\,\mathrm{d}\log\frac{\mathrm{M}(a,c)}{\mathrm{M}(a,b)}\end{aligned}
\]
\Needspace{21\baselineskip}
The other relation also extends to the differentials of the logarithms of products of two terms. It follows from $b'b'=ab$ and ${}'c{}'c=4ac$ in the form
\Needspace{10\baselineskip}
\[
\begin{aligned}
\mathrm{d}\log(a^{\mathrm{m}}\mathbin{.} b^{\mathrm{m}})&=\mathrm{d}\log(a^{\mathrm{m}+1}\mathbin{.} b^{\mathrm{m}+1})\\&\quad-\Delta\mathbin{.}2^{m+1}\mathbin{.} c^{\mathrm{m}+1}\mathbin{.} c^{\mathrm{m}+1}\\[3pt]
\mathrm{d}\log({}^{\mathrm{m}}a\mathbin{.}{}^{\mathrm{m}}c)&=\mathrm{d}\log({}^{\mathrm{m}+1}a\mathbin{.}{}^{\mathrm{m}+1}c)\\&\quad+\Delta\mathbin{.}\frac{1}{2^{m+1}}\mathbin{.}{}^{\mathrm{m}+1}b\mathbin{.}{}^{\mathrm{m}+1}b
\end{aligned}
\]
\Needspace{20\baselineskip}
and, when $m$ is allowed to grow to the infinite limit, gives the equation stated in Article 9 and the corresponding one, namely:
\Needspace{10\baselineskip}
\[
\begin{aligned}
2\,\mathrm{d}\log\mathrm{M}(a,b)&=\mathrm{d}\log(a^{\mathrm{n}}\mathbin{.} b^{\mathrm{n}})\\&\quad+\Delta\{2^{n+1}\mathbin{.} c^{\mathrm{n}+1}\mathbin{.} c^{\mathrm{n}+1}+\ .\ +2^{n+m}\mathbin{.} c^{\mathrm{n}+\mathrm{m}}\mathbin{.} c^{\mathrm{n}+\mathrm{m}}+\ .\}\\[3pt]
2\,\mathrm{d}\log\mathrm{M}(a,c)&=\mathrm{d}\log({}^{\mathrm{n}}a\mathbin{.}{}^{\mathrm{n}}c)\\&\quad-\Delta\{2^{-n-1}\mathbin{.}{}^{\mathrm{n}+1}b\mathbin{.}{}^{\mathrm{n}+1}b+\ .\ +2^{-n-m}\mathbin{.}{}^{\mathrm{n}+\mathrm{m}}b\mathbin{.}{}^{\mathrm{n}+\mathrm{m}}b+\ .\}
\end{aligned}
\]
\Needspace{22\baselineskip}
\par From this, by eliminating the differentials with the aid of the expression for $\Delta$ found above, the equation
\Needspace{10\baselineskip}
\[
\begin{aligned}
\frac{4}{\pi}\mathrm{M}(a,b)\,\mathrm{M}(a,c)&=.\ .-2^{-\mu+n}\mathbin{.}{}^{\mu-\mathrm{n}}b\mathbin{.}{}^{\mu-\mathrm{n}}b\\&\quad-.\ .-2^{n-1}\mathbin{.} b^{\mathrm{n}-1}\mathbin{.} b^{\mathrm{n}-1}\\
&{}+2^{n}\mathbin{.} a^{\mathrm{n}}\mathbin{.} a^{\mathrm{n}}-2^{n+1}\mathbin{.} c^{\mathrm{n}+1}\mathbin{.} c^{\mathrm{n}+1}\\&\quad-.\ .-2^{n+m}\mathbin{.} c^{\mathrm{n}+\mathrm{m}}\mathbin{.} c^{\mathrm{n}+\mathrm{m}}-.\ .
\end{aligned}
\]
can be derived. \textsc{Gauss} wrote it down alongside the first equation for $\frac{\pi}{2}\frac{\mathrm{M}(a,b)}{\mathrm{M}(a,c)}$ reproduced in the preceding article, without indicating the route by which they had been found.]
\Needspace{30\baselineskip}
\AGMArticle{14}
\par [For the total differentials of the second order, it follows directly from the expressions for $\Delta$ that every system of terms $a,b,c$ bearing a common index satisfies the three equations
\Needspace{10\baselineskip}
\[
\begin{aligned}
\frac{1}{\Delta}\,\mathrm{d}\log\frac{c}{b}\mathbin{.}\mathrm{d}\log a&=\tfrac12\,\mathrm{d}\left(\frac{1}{\Delta}\,\mathrm{d}\log\frac{c}{b}\right)\\[4pt]
\frac{1}{\Delta}\,\mathrm{d}\log\frac{c}{a}\mathbin{.}\mathrm{d}\log b&=\tfrac12\,\mathrm{d}\left(\frac{1}{\Delta}\,\mathrm{d}\log\frac{c}{a}\right)\\[4pt]
\frac{1}{\Delta}\,\mathrm{d}\log\frac{a}{b}\mathbin{.}\mathrm{d}\log c&=\tfrac12\,\mathrm{d}\left(\frac{1}{\Delta}\,\mathrm{d}\log\frac{a}{b}\right)
\end{aligned}
\]
\Needspace{15\baselineskip}
If the numerators and denominators under the logarithms to be differentiated twice are multiplied in turn by $a,b,c$, and all logarith\SourceLeaf{390}ms of quotients are split into differences of logarithms, it is easily seen that the expression
\Needspace{10\baselineskip}
\[
\frac{1}{\Delta}\,\mathrm d\log a^{\mathrm n}\mathbin{.}\mathrm d\log b^{\mathrm n}
-\frac12\mathrm d\left(\frac{1}{\Delta}\,\mathrm d\log a^{\mathrm n}b^{\mathrm n}\right)
\]
which is to be denoted by $D$, does not change its value if
$a^{\mathrm n}$ and $b^{\mathrm n}$ are replaced by $a^{\mathrm n}$ and $c^{\mathrm n}$ or by
$b^{\mathrm n}$ and $c^{\mathrm n}$. Taking the mean of the last two expressions so obtained and noting that
\Needspace{10\baselineskip}
\[
\mathrm d\log a^{\mathrm n}+\mathrm d\log b^{\mathrm n}
=2\mathrm d\log b^{\mathrm n+1}
\]
holds, it follows that the two quantities named may also be replaced in that expression, without changing its value, by $c^{\mathrm n}$ and $b^{\mathrm n+1}$; and further, if $\mathrm d\log c^{\mathrm n}$ is again replaced in this by
$\frac12\mathrm d\log a^{\mathrm n+1}+\frac12\mathrm d\log c^{\mathrm n+1}$,
that $a^{\mathrm n+1}$ and $b^{\mathrm n+1}$ can be substituted with the same result for
$a^{\mathrm n}$ and $b^{\mathrm n}$. \Needspace{23\baselineskip}
The value of that differential expression is therefore independent of $n$ and hence equal to
\Needspace{23\baselineskip}
\[
\begin{aligned}
D&=\displaystyle\frac{1}{\Delta}\,\mathrm d\log b^{\mathrm n}\mathbin{.}\mathrm d\log c^{\mathrm n}
-\frac12\mathrm d\left(\frac{1}{\Delta}\,\mathrm d\log b^{\mathrm n}c^{\mathrm n}\right)\\[6pt]
&=\displaystyle\frac{1}{\Delta}\,\mathrm d\log c^{\mathrm n}\mathbin{.}\mathrm d\log a^{\mathrm n}
-\frac12\mathrm d\left(\frac{1}{\Delta}\,\mathrm d\log c^{\mathrm n}a^{\mathrm n}\right)\\[6pt]
&=\displaystyle\frac{1}{\Delta}\,\mathrm d\log a^{\mathrm n}\mathbin{.}\mathrm d\log b^{\mathrm n}
-\frac12\mathrm d\left(\frac{1}{\Delta}\,\mathrm d\log a^{\mathrm n}b^{\mathrm n}\right)\\[6pt]
&=\displaystyle\mathrm M(a,b)\,\mathrm d\left\{\frac{1}{\Delta}\,\mathrm d\frac{1}{\mathrm M(a,b)}\right\}\\&=\mathrm M(a,c)\,\mathrm d\left\{\frac{1}{\Delta}\,\mathrm d\frac{1}{\mathrm M(a,c)}\right\}
\end{aligned}
\]
\Needspace{22\baselineskip}
\par If $a$ is assumed constant, then
\Needspace{10\baselineskip}
\[
D=\Delta bbcc=-bb\,\mathrm d\log b=cc\,\mathrm d\log c
\]
and the differential equation derived in Article 8 from the series expansion is obtained, with
$\frac{a}{\mathrm M(a,c)}$ and $\frac{a}{\mathrm M(a,b)}$ as values of $\mu$:
\Needspace{10\baselineskip}
\[
\mathrm{dd}\mu-\frac{\mathrm d\Delta}{\Delta}\mathbin{.}\mathrm d\mu
-\Delta\Delta bbcc\mathbin{.}\mu=0
\]
From this, by the investigations in the paper
`\textit{Determinatio seriei nostrae per aequationem differentialem secundi ordinis}',
one again obtains the representations by the \textsc{Gauss}ian series:
\Needspace{10\baselineskip}
\[
\begin{gathered}\frac{a}{\mathrm M(a,b)}=F\left(\frac12,\frac12,1,\frac{cc}{aa}\right),\\[7pt]
\frac{a}{\mathrm M(a,c)}=F\left(\frac12,\frac12,1,\frac{bb}{aa}\right)\end{gathered}
\]
\SourceLeaf{391}
and, as special cases of equations [90] and [96] there, the limit found above in Article 12 of $\mathrm M(1,\varepsilon)\log\frac{4}{\varepsilon}$ and the first-order differential equation between
$\mathrm M(a,b)$ and $\mathrm M(a,c)$ stated in the preceding article.]

\clearpage
\endgroup
